\chapter{Market making automatisé sur options}\label{ch:hf:automated-options-market-making}

Un teneur de marché d'options cote chaque prix d'exercice et chaque échéance d'une action à la fois, plusieurs milliers de cours sur une douzaine de bourses, et un balayage peut en toucher quarante en une même milliseconde~; la protection de la bourse est ce qui arrête la quarante et unième.
Dans la chaîne de 170 séries de ce chapitre, un trader qui
sait que l'action a bougé de 1~\% et frappe chaque cotation périmée, une toutes les cinq microsecondes, prendrait 79\,000~\$ à un teneur de marché sans protection~;
une protection qui retire toutes les cotations après treize exécutions limite la perte à 12\,900~\$ et, sur une journée ordinaire d'ordres clients multi-jambes, ne coûte
que trois exécutions sur 8\,235.

\section{La surface en temps réel}

Les valeurs théoriques d'un teneur de marché d'options proviennent d'une surface de volatilité (le livre~5, chapitres~7 et~8)~: un smile paramétré par
échéance, ajusté aux cotations et aux transactions du marché, déplacé avec le sous-jacent selon une règle (sticky strike, sticky delta) et réajusté lorsque
le marché montre que la règle était fausse.

\begin{definition}[Ajustement de la surface en temps réel]\label{def:hf:automated-options-market-making:live}
Un \emph{ajustement de la surface en temps réel}\index{ajustement de la surface en temps réel} (live surface fit) est la ré-estimation continue de la surface de volatilité d'un teneur de marché d'options pendant la
journée de négociation, mettant à jour les paramètres de chaque échéance à partir des cotations et des transactions les plus récentes à mesure qu'elles arrivent, assez vite pour que chaque cotation
envoyée par le teneur de marché soit valorisée sur une surface cohérente avec le marché courant.
\end{definition}

À haute fréquence, le travail est scindé~: le prix du sous-jacent déplace toute la surface selon la règle, microseconde après microseconde, et la
forme du smile est réajustée à un rythme plus lent, typiquement lorsque suffisamment d'informations nouvelles sont arrivées. Une tranche SVI brute ajustée par la méthode quasi-explicite
(\texttt{firm.svi} du livre~5) à dix-sept volatilités cotées, chacune avec un bruit de 0.2 point de volatilité, retrouve, à 0.025 point près, le vrai smile~:
l'ajustement moyenne le bruit que portent les cotations individuelles. Muravyev et Pearson ont montré pourquoi la fair value compte tant pour les preneurs de liquidité~: les prix des options
évoluent de façon prévisible à haute fréquence~: de nombreux traders choisissent le moment de leurs exécutions pour acheter quand la fair value est proche du cours vendeur et vendre quand elle est proche du
cours acheteur, et les fourchettes effectives de ceux qui choisissent le moment de leurs exécutions sont inférieures à 40~\% des mesures conventionnelles.

\section{Cotation de masse et protection des cotations}

\begin{definition}[Cotation de masse]\label{def:hf:automated-options-market-making:mass}
Une \emph{cotation de masse}\index{cotation de masse} (mass quote) est un message unique par lequel un teneur de marché envoie ou remplace des cours acheteurs et vendeurs sur de nombreuses séries d'options à
la fois, que la bourse traite comme une unité, de sorte qu'un teneur de marché puisse mettre à jour des milliers de cotations après un mouvement du sous-jacent sans envoyer
des milliers d'ordres distincts.
\end{definition}

Une chaîne de 17 prix d'exercice et cinq échéances, calls et puts, compte 170 séries~; cotée à un demi-point de volatilité de part et d'autre de la valeur
théorique (\cref{lst:hf:automated-options-market-making:quote}), le call à la monnaie à 30 jours, à 2.92~\$, a une largeur de 11~cents, et les largeurs vont du
plancher d'un cent pour les options profondes à 20~cents. Chaque cotation reste sur la bourse jusqu'à la \omterm{def:hf:automated-options-market-making:mass}{cotation de masse} suivante qui la remplace~; entre le mouvement du sous-jacent et ce
remplacement, les 170 sont périmées en même temps.

La réponse de la bourse est la protection~: un ensemble de compteurs par teneur de marché et par classe (nombre d'exécutions, de contrats, notionnel, ou pourcentage
de la quantité cotée exécutée) sur une courte fenêtre, et une règle qui annule toutes ses cotations de la classe lorsqu'un compteur atteint sa limite.

\begin{dated}{2026-09}{Paramètres de risque d'une bourse}\label{dat:hf:automated-options-market-making:cboe}
La page de Cboe consacrée aux outils de gestion des risques de ses bourses d'options recense quatre types de paramètres~: le notionnel (contrats négociés multipliés par la prime), le volume (contrats
négociés), le compte (nombre d'exécutions) et le pourcentage de cotation (la somme, sur toutes les séries d'une même racine, du pourcentage de contrats cotés
exécutés), chacun sur une période fixée entre 100 millisecondes et une journée entière~; dès qu'un paramètre est atteint, aucune nouvelle transaction n'est exécutée et toute
transaction en route est rejetée.
\end{dated}

Le balayage du chapitre (\cref{lst:hf:automated-options-market-making:sweep})~: l'action bouge de 1~\%~; 150 des 170 cotations sont désormais mal valorisées~; un trader
informé les frappe, les plus rentables d'abord, une toutes les cinq microsecondes~; la bourse retire toutes les cotations après $n$ exécutions dans les 100 millisecondes.
La perte croît presque linéairement en $n$, d'environ 1\,000~\$ par exécution, parce que les cibles les plus rentables sont des options profondément dans la monnaie qui bougent
dollar pour dollar avec l'action. Sur une journée ordinaire, les ordres clients arrivent en moyenne toutes les cinq secondes, 70~\% d'entre eux portant sur une seule série et les autres
pour des spreads et des strips à deux, quatre et douze jambes exécutés en moins de deux millisecondes~; chaque exécution rapporte 4~\$. Un seuil inférieur à douze retire les
cotations des ordres légitimes à douze jambes et laisse le teneur de marché hors circuit pendant une demi-seconde à chaque fois
(\cref{fig:hf:automated-options-market-making:protection}).

\begin{omfigure}
\begin{tikzpicture}
\begin{axis}[width=11cm,height=5.4cm,axis y line*=right,axis x line=none,xmode=log,log basis x=10,xmin=1.6,xmax=50,ymin=0,ymax=2600,
  ylabel={\small exécutions perdues un jour ordinaire},tick label style={font=\footnotesize},table/col sep=comma]
\addplot[omDef,thick,dashed,mark=square*,mark size=1.3pt] table[x=n,y=lost]{figdata/hft/19-automated-options-market-making/protection.csv};
\label{plot:hf:protlost}
\end{axis}
\begin{axis}[width=11cm,height=5.4cm,axis y line*=left,xmode=log,log basis x=10,xmin=1.6,xmax=50,xtick={2,3,5,10,13,20,40},
  xticklabels={2,3,5,10,13,20,40},xlabel={\small seuil de protection (exécutions en 100 ms)},ylabel={\small perte du balayage (milliers de \$)},ymin=0,
  ymax=42,tick label style={font=\footnotesize},grid=major,grid style={gray!25},table/col sep=comma,
  legend style={font=\scriptsize,at={(0.5,-0.3)},anchor=north},legend columns=2]
\addplot[omThm,thick,mark=*,mark size=1.3pt] table[x=n,y=loss]{figdata/hft/19-automated-options-market-making/protection.csv};
\addlegendentry{perte d'un balayage de 1~\% (gauche)}
\addlegendimage{/pgfplots/refstyle=plot:hf:protlost}\addlegendentry{exécutions perdues un jour ordinaire (droite)}
\end{axis}
\end{tikzpicture}

\omcaption{La perte du teneur de marché face à un balayage informé après un mouvement de 1~\% de l'action (axe de gauche, milliers de dollars) et les exécutions clients qu'il perd
un jour ordinaire pendant que ses cotations sont retirées (axe de droite), face au seuil de comptage de la bourse~; 170 séries cotées chacune à dix contrats~; le
balayage frappe une série toutes les cinq microsecondes~; ordres ordinaires d'une à douze jambes, 8\,235 exécutions par jour sans protection. Sans protection, le
balayage prend 79\,370~\$. Données~: \texttt{hf\_options.protection\_table}.}\label{fig:hf:automated-options-market-making:protection}
\end{omfigure}

\section{Couverture en delta à grande vitesse}

Le balayage laisse le teneur de marché à découvert d'environ 13\,000 actions de delta pour un seuil de treize (il a vendu des calls et acheté des puts après
que l'action a monté). La perte est déjà faite~: les options ont été vendues à des prix inférieurs à la nouvelle valeur. Le rôle de la couverture est d'empêcher le mouvement suivant d'y ajouter, et à
grande vitesse, la couverture part dès l'arrivée des exécutions, sur le sous-jacent ou sur l'instrument corrélé le plus liquide, dans la bande de couverture
décrite au livre~5, chapitre~26. Muravyev a montré que l'inventaire pèse davantage sur les prix des options qu'on ne le pensait~: en décomposant l'impact des
transactions d'options entre risque d'inventaire et information asymétrique, il trouve les deux importants et la composante d'inventaire plus grande, et les déséquilibres des ordres dus au
risque d'inventaire font bouger les prix des options cinq fois plus que les estimations antérieures.

\section{Risque par prix d'exercice et par échéance}

\begin{definition}[Limite de risque par prix d'exercice]\label{def:hf:automated-options-market-making:strike}
Une \emph{limite de risque par prix d'exercice}\index{limite de risque par prix d'exercice} (strike risk limit) est le plafond d'exposition (vega, gamma ou nombre de contrats) que le teneur de marché accepte d'accumuler sur un
prix d'exercice ou une échéance d'une chaîne~; lorsqu'une exécution la franchit, le teneur de marché annule ou élargit ses cotations sur ce prix d'exercice ou cette échéance jusqu'à ce que
l'exposition soit réduite.
\end{definition}

La protection fondée sur le comptage compte des exécutions, non le risque. Un balayage de volatilité, après une hausse de deux points de volatilité implicite, frappe 130 cotations de chaque
échéance et laisse le teneur de marché à découvert de 10\,600~\$ de vega par point de volatilité, pour une perte de 16\,500~\$. Une limite de 1\,000~\$ de vega par
échéance, imposée par son propre système 50 microsecondes après l'exécution qui la franchit, ramène la perte à 9\,900~\$~; imposée en 5 microsecondes, à
6\,700~\$. La valeur de la limite dépend de sa vitesse~: chaque microseconde de réaction est une cotation périmée frappée de plus.

\section{Ce que les grandes firmes ont publié}

Les teneurs de marché d'options publient peu sur leurs systèmes de cotation~; ce qui est public, ce sont les règles et les spécifications des bourses (les protections de
l'encadré daté, les messages de \omterm{def:hf:automated-options-market-making:mass}{cotation de masse}), les dépôts auprès des régulateurs et les travaux académiques sur les données que ces firmes produisent. Le chapitre ne cite que ceux-là.

\section{Fiches de stratégie}

\begin{strategyfile}{Cotation sur surface ajustée d'une chaîne cotée}\label{strat:hf:automated-options-market-making:surface}
\sfield{Qui vous paie, et pourquoi} Les utilisateurs d'options (couvreurs, traders directionnels, traders de volatilité) qui ont besoin d'un prix sur n'importe quelle série, à tout moment.
\sfield{Instruments et plateformes} La chaîne cotée d'une action ou d'un indice sur chaque bourse d'options.
\sfield{Signal} Les valeurs théoriques issues d'une surface en temps réel~; le mouvement du sous-jacent selon la règle de la surface.
\sfield{Dimensionnement et exécution} \omterm{def:hf:automated-options-market-making:mass}{Cotations de masse} à une largeur fixée par le vega et par l'incertitude de la surface~; recalibrer le smile à un rythme plus lent.
\sfield{Coûts} Frais de bourse, largeur abandonnée aux concurrents, sélection adverse de la part des traders qui calibrent le moment de leurs exécutions.
\sfield{Comment elle meurt} Cotations périmées après les mouvements~; concurrents dont les recalibrages sont plus rapides.
\sfield{Horizon, capacité, infrastructure} De la microseconde à la journée~; des milliers de séries~; données de marché à l'échelle de l'OPRA.
\sfield{Backtest honnête} La surface telle qu'elle était connue à chaque cotation~; les exécutions là où les traders qui jouent le timing les auraient prises.
\sfield{Sources} Muravyev et Pearson (2020)~; le livre~5, chapitres 7, 8 et 26.
\end{strategyfile}

\begin{strategyfile}{Cotation de masse avec protection des cotations}\label{strat:hf:automated-options-market-making:mass}
\sfield{Qui vous paie, et pourquoi} Comme ci-dessus, la protection de la bourse limitant la perte aux balayages.
\sfield{Instruments et plateformes} Bourses avec messages de \omterm{def:hf:automated-options-market-making:mass}{cotation de masse} et paramètres de risque.
\sfield{Signal} Les mêmes valeurs théoriques~; les seuils de la protection fixés à partir du plus gros ordre légitime.
\sfield{Dimensionnement et exécution} Fixer le seuil de comptage juste au-dessus du plus gros ordre multi-jambes (treize contre des ordres à douze jambes dans le modèle).
\sfield{Coûts} Exécutions perdues lorsque la protection se déclenche sur les flux légitimes (326 par jour au seuil de dix, trois au seuil de treize).
\sfield{Comment elle meurt} Balayages plus rapides que la première exécution de la protection, et protections fixées trop larges.
\sfield{Horizon, capacité, infrastructure} Millisecondes~; une protection par classe et par bourse.
\sfield{Backtest honnête} La logique de protection de la bourse rejouée sur les exécutions de la journée, y compris les ordres multi-jambes.
\sfield{Sources} L'encadré daté~; le présent chapitre.
\end{strategyfile}

\begin{strategyfile}{Livre d'options couvert en delta}\label{strat:hf:automated-options-market-making:delta}
\sfield{Qui vous paie, et pourquoi} Les clients d'options, par la fourchette~; la couverture empêche la valeur du livre de suivre l'action.
\sfield{Instruments et plateformes} Le livre d'options et le sous-jacent (ou son contrat à terme).
\sfield{Signal} Le delta du livre après chaque exécution.
\sfield{Dimensionnement et exécution} Couvrir immédiatement au-delà d'une bande, dans l'instrument le moins cher~; agréger d'abord sur toute la chaîne.
\sfield{Coûts} La fourchette et les frais du sous-jacent~; le gamma entre les couvertures.
\sfield{Comment elle meurt} Exécutions corrélées qui font bouger le delta du livre plus vite que la couverture ne peut suivre.
\sfield{Horizon, capacité, infrastructure} De la microseconde à la seconde~; la saisie d'ordres du sous-jacent à côté de celle des options.
\sfield{Backtest honnête} Prix de couverture après l'horodatage de l'exécution, plus la latence de couverture.
\sfield{Sources} Muravyev (2016)~; le livre~5, chapitre 26.
\end{strategyfile}

\begin{strategyfile}{Cotation d'options tenant compte des événements}\label{strat:hf:automated-options-market-making:event}
\sfield{Qui vous paie, et pourquoi} Les traders qui se positionnent sur des résultats ou des événements macroéconomiques, et qui paient pour le mouvement implicite.
\sfield{Instruments et plateformes} Échéances encadrant un événement connu.
\sfield{Signal} La variance d'événement dans la structure par terme (l'arithmétique de la variance d'événement du livre~5), à comparer à l'historique des mouvements de l'action.
\sfield{Dimensionnement et exécution} Coter les échéances d'événement sur une surface avec la variance de l'événement~; resserrer les limites de vega avant l'événement~; retirer les cotations ou les élargir à
l'annonce (chapitre~18).
\sfield{Coûts} Se tromper sur l'ampleur de l'événement~; le gamma à cheval sur l'annonce.
\sfield{Comment elle meurt} Surprises supérieures au mouvement implicite.
\sfield{Horizon, capacité, infrastructure} Des jours jusqu'à l'événement~; des secondes autour de celui-ci.
\sfield{Backtest honnête} Le calendrier des événements tel qu'il était connu à l'époque.
\sfield{Sources} Chapitre~18~; le livre~5, chapitre~8.
\end{strategyfile}

\begin{strategyfile}{Cotation du skew limitée par le vega}\label{strat:hf:automated-options-market-making:vega}
\sfield{Qui vous paie, et pourquoi} Les acheteurs de protection (puts) et les vendeurs de calls qui maintiennent le skew à un niveau élevé.
\sfield{Instruments et plateformes} Puts et calls hors de la monnaie sur plusieurs échéances.
\sfield{Signal} Le niveau du skew comparé à son historique et à l'ajustement de la surface.
\sfield{Dimensionnement et exécution} Coter les ailes avec des limites de vega par prix d'exercice et par échéance, appliquées en quelques microsecondes.
\sfield{Coûts} Risque de queue sur les ailes~; couvertures dans le sous-jacent et sur d'autres échéances.
\sfield{Comment elle meurt} Un balayage de volatilité~: 16\,500~\$ non protégés dans le modèle~; la vitesse des limites détermine la perte (9\,900~\$ à 50 microsecondes,
6\,700~\$ à 5).
\sfield{Horizon, capacité, infrastructure} Des jours~; application rapide des limites.
\sfield{Backtest honnête} Limites appliquées avec leur temps de réaction réel.
\sfield{Sources} Le présent chapitre.
\end{strategyfile}

\section{Tutoriel~: quarante prix d'exercice en une milliseconde}\label{tut:hf:automated-options-market-making:tut}

\begin{tutorial}
\textbf{Objectif.} Coter une chaîne à partir d'une surface, puis mesurer ce qu'un balayage informé prend face à chaque seuil de protection et ce que le seuil
coûte les jours ordinaires. \textbf{État final~:} la \cref{fig:hf:automated-options-market-making:protection} et les chiffres du texte.
\begin{tutsteps}
\item \textbf{La \omterm{def:hf:automated-options-market-making:mass}{cotation de masse}} à partir des valeurs théoriques du livre~5.
\omcode{code/firm/optquoter/firm_optquoter.py}{63}{74}{Valeur théorique et vega à partir de \texttt{firm.optmm.theo}, une demi-largeur en points de volatilité multipliée par le vega, un plancher d'un centime.}\label{lst:hf:automated-options-market-making:quote}
\item \textbf{Le balayage} face au nombre de la bourse et aux limites de vega du teneur de marché.
\omcode{code/firm/optquoter/firm_optquoter.py}{108}{126}{Les cotations périmées sont touchées en premier par le profit le plus élevé~; le nombre de la bourse et la limite de vega par échéance l'arrêtent.}\label{lst:hf:automated-options-market-making:sweep}
\item \textbf{Une journée ordinaire}~: ordres clients multi-jambes face aux mêmes seuils (\texttt{hf\_options.protection\_table}).
\item \textbf{Le balayage de volatilité} avec des limites de vega à deux vitesses de réaction (\texttt{hf\_options.vol\_sweep})~; le recalibrage du smile
(\texttt{hf\_options.refit}).
\end{tutsteps}
\textbf{Pour aller plus loin.} Remplacer le comptage par un paramètre en pourcentage de la cotation~; faire arriver le balayage sur plusieurs bourses à la fois, chacune avec
sa propre protection~; faire tourner la chaîne sur \texttt{firm.exchsim} avec son message de \omterm{def:hf:automated-options-market-making:mass}{cotation de masse}.
\end{tutorial}

\begin{center}
\small
\begin{tabular}{@{}lrrrrrrr@{}}
\toprule
seuil (exécutions) & 2 & 5 & 10 & 13 & 20 & 40 & aucun\\
\midrule
perte d'un balayage de 1~\% (\$) & 1\,980 & 4\,950 & 9\,899 & 12\,867 & 19\,773 & 38\,718 & 79\,370\\
exécutions perdues un jour ordinaire & 2\,441 & 1\,052 & 326 & 3 & 0 & 0 & 0\\
fourchette perdue par jour (\$) & 9\,764 & 4\,208 & 1\,304 & 12 & 0 & 0 & 0\\
\bottomrule
\end{tabular}
\end{center}

\section{Construction~: le cotateur d'options}\label{bld:hf:automated-options-market-making:optquoter}

\begin{build}
\bfield{Rôle} Coter une chaîne à partir d'une surface en temps réel en \omterm{def:hf:automated-options-market-making:mass}{cotations de masse}, et mesurer les protections et les limites face aux balayages et aux flux ordinaires.
\bfield{Interface} \texttt{Chain(spot, strikes, expiries\_days)}, \texttt{surface(atm, term, skew, ref)},
\texttt{mass\_quote(chain, spot, surf, half\_width\_vol, size)}, \texttt{refit\_slice(strikes, spot, tau, vols)}, \texttt{Protection(count, window\_ms)},
\texttt{sweep(quotes, jump, chain, surf, protection, gap\_us, size, dvol, vega\_limit, react\_us)}, \texttt{normal\_day(n\_series, protection, seed)}, \texttt{bucket\_vega}.
Construit sur \texttt{firm.optmm}, \texttt{firm.volpnl}, \texttt{firm.svi} et \texttt{firm.bs} du livre~5.
\bfield{Règles} Prix par action, 100 actions par contrat~; le teneur de marché est la contrepartie de chaque exécution d'un balayage.
\bfield{Tests d'acceptation} \texttt{code/firm/optquoter/tests/}~: les cotations encadrent la valeur théorique, avec des largeurs proportionnelles au vega~; le balayage s'arrête au nombre prévu, et
sa perte égale la somme des fourchettes exécutées~; un seuil plus élevé ne perd jamais moins~; une limite de vega qui réagit plus vite perd moins~; la journée ordinaire ne perd aucune exécution
lorsque le seuil dépasse le nombre de jambes de chaque ordre~; le recalibrage SVI retrouve une tranche à l'intérieur de son bruit.
\bfield{Pour aller plus loin} Paramètres en pourcentage de la cotation et en notionnel~; plusieurs bourses~; recalibrage incrémental de la surface à chaque transaction.
\end{build}

\begin{omsources}
\item D.~Muravyev, Order flow and expected option returns, \emph{Journal of Finance} 71(2), 2016, 673--708.
\item D.~Muravyev, N.~D.~Pearson, Options trading costs are lower than you think, \emph{Review of Financial Studies} 33(11), 2020, 4973--5014.
\item Cboe, US Options risk management tools (page web, consulté le 2026).
\end{omsources}

\section{Exercices}

\begin{exercise}[$\star$]\label{exo:hf:automated-options-market-making:1}
Le call à la monnaie à 30 jours a un vega de 0.114 par point de volatilité. Quelle est sa largeur à un demi-point de part et d'autre~?
\end{exercise}

\begin{exercise}[$\star$]\label{exo:hf:automated-options-market-making:2}
Un call profondément dans la monnaie, de delta égal à un, est coté 20.01~\$ au cours vendeur, pour dix contrats. Le titre monte de 1~\$. Que gagne un sweeper sur cette position~?
\end{exercise}

\begin{exercise}[$\star$]\label{exo:hf:automated-options-market-making:3}
Pourquoi le seuil de comptage doit-il dépasser le plus grand ordre multi-jambes légitime~?
\end{exercise}

\begin{exercise}[$\star\star$]\label{exo:hf:automated-options-market-making:4}
Pourquoi la perte du balayage croît-elle presque linéairement avec le seuil~?
\end{exercise}

\begin{exercise}[$\star\star$]\label{exo:hf:automated-options-market-making:5}
Pourquoi une limite de vega par échéance, appliquée en 50~microsecondes, perd-elle davantage qu'une limite appliquée en 5~microsecondes~?
\end{exercise}

\begin{exercise}[$\star\star$]\label{exo:hf:automated-options-market-making:6}
Muravyev et Pearson ont établi que les exécutions au moment choisi coûtent moins de 40~\% des spreads conventionnels. Que cela implique-t-il pour les largeurs d'un teneur de marché~?
\end{exercise}

\begin{exercise}[$\star\star\star$]\label{exo:hf:automated-options-market-making:7}
\emph{Programmation.} Relancer le balayage avec un mouvement de 2~\%. Au seuil de 13, quelle est la perte, et sans protection~?
\end{exercise}

\begin{exercise}[$\star\star\star$]\label{exo:hf:automated-options-market-making:8}
\emph{Trouver la faille.} «~Notre seuil de protection est de quarante exécutions~; nous n'avons jamais perdu plus de 5\,000~\$ à cause d'un balayage, donc il est sûr.~»
\end{exercise}

\section{Problème~: quarante prix d'exercice en une milliseconde}

\begin{problem}[{Problème du week-end --- quarante prix d'exercice en une milliseconde}]\label{pb:hf:automated-options-market-making:1}
Un teneur de marché d'options fixe ses protections et ses limites pour une chaîne de 170 séries.

\medskip\noindent\textbf{Partie I --- Cotations.}
\begin{enumerate}
\item Définir un \omterm{def:hf:automated-options-market-making:live}{ajustement de la surface en temps réel} et la façon dont son travail se répartit entre les horloges.
\item Dans quelle mesure un recalibrage SVI retrouve-t-il une tranche bruitée~?
\item Définir une \omterm{def:hf:automated-options-market-making:mass}{cotation de masse}.
\item Donner les largeurs de la chaîne.
\end{enumerate}

\medskip\noindent\textbf{Partie II --- Protection.}
\begin{enumerate}[resume]
\item Résumer l'encadré daté.
\item Décrire le balayage et la journée ordinaire.
\item Donner la perte due au balayage et les exécutions perdues aux seuils de 5, 10, 13 et 20.
\item Pourquoi les cibles les plus rentables viennent-elles en premier~?
\end{enumerate}

\medskip\noindent\textbf{Partie III --- Risque.}
\begin{enumerate}[resume]
\item Quel delta le balayage laisse-t-il, et que fait la couverture~?
\item Qu'ont établi Muravyev et Pearson sur l'inventaire et les prix des options~?
\item Définir une \omterm{def:hf:automated-options-market-making:strike}{limite de risque par prix d'exercice}.
\item Donner les pertes du balayage de volatilité avec des limites à deux vitesses.
\end{enumerate}

\medskip\noindent\textbf{Partie IV --- Le verdict.}
\begin{enumerate}[resume]
\item Énoncer le \emph{résultat nommé}~: la perte d'un balayage d'information planté en fonction du seuil de protection des cotations, et les exécutions abandonnées par
un seuil serré les jours ordinaires.
\item Quel seuil choisiriez-vous, et pourquoi~?
\item Que changeraient plusieurs bourses, chacune avec sa propre protection~?
\item Pourquoi les propres limites du teneur de marché sont-elles nécessaires en plus de celles de la bourse~?
\item Qu'ont établi Muravyev et Pearson sur la prise de liquidité dans les options~?
\item Quelle fiche de stratégie est la plus exposée aux balayages~?
\item Comment fixeriez-vous les limites de vega par échéance~?
\item En une phrase~: que fait arbitrer la protection des cotations~?
\end{enumerate}
\end{problem}

\section{Questions d'entretien}

\begin{interviewq}[$\star$ \iqroles{trader}]\label{iq:hf:automated-options-market-making:1}
Le titre bondit de 3~\% en une milliseconde. Que se passe-t-il, dans l'ordre, pour vos cotations d'options~?
\end{interviewq}

\begin{interviewq}[$\star\star$ \iqroles{développeur}]\label{iq:hf:automated-options-market-making:2}
Concevez le chemin qui va d'un tick du sous-jacent à une \omterm{def:hf:automated-options-market-making:mass}{cotation de masse} de 5\,000 séries. Où part le temps~?
\end{interviewq}

\begin{interviewq}[$\star\star$ \iqroles{chercheur}]\label{iq:hf:automated-options-market-making:3}
Comment décideriez-vous de recalibrer un smile plutôt que de le déplacer avec le sous-jacent~?
\end{interviewq}

\begin{interviewq}[$\star\star$ \iqroles{risque}]\label{iq:hf:automated-options-market-making:4}
Votre protection s'est déclenchée quarante fois aujourd'hui sur une même classe. Que examinez-vous~?
\end{interviewq}

\begin{interviewq}[$\star\star$ \iqroles{trader}]\label{iq:hf:automated-options-market-making:5}
Pourquoi pourriez-vous coter les puts plus larges que les calls au même vega~?
\end{interviewq}

\begin{interviewq}[$\star\star\star$ \iqroles{chercheur}]\label{iq:hf:automated-options-market-making:6}
Dérivez la perte espérée d'un balayage contre un seuil de comptage $n$ lorsque les cibles arrivent par ordre de edge décroissant $e_1\ge e_2\ge\cdots$.
\end{interviewq}
